\chapter{Gesamtinhaltsverzeichnis}

Die folgenden Verzeichnisse geben den Inhalt aller achtundzwanzig Bände wieder. Die Seitenzahlen beziehen sich jeweils auf den einzelnen Band.

\section*{Band 0 --- Die Kunst des Rechnens}
\addcontentsline{toc}{section}{Band 0 --- Die Kunst des Rechnens}
\nopagebreak
{\small\itshape\color{grau} Vorkurs Schulmathematik $\cdot$ 154 Seiten}\par\vspace{4pt}
\nopagebreak
{\small\renewcommand{\arraystretch}{1.28}
\begin{longtable}{@{}P{0.9cm}P{6.6cm}P{1.0cm}@{}}
\multicolumn{3}{@{}l}{\color{akzent}\bfseries\small I.~Zahlen und Rechnen}\\[1pt]
1 & Die Zahlbereiche & 7\\
2 & Bruchrechnung & 15\\
3 & Potenzen, Wurzeln, Logarithmen & 23\\
\multicolumn{3}{@{}l}{\color{akzent}\bfseries\small II.~Algebra}\\[1pt]
4 & Terme umformen & 33\\
5 & Gleichungen & 39\\
6 & Ungleichungen und Systeme & 47\\
\multicolumn{3}{@{}l}{\color{akzent}\bfseries\small III.~Funktionen}\\[1pt]
7 & Der Funktionsbegriff & 57\\
8 & Die wichtigsten Funktionstypen & 67\\
\multicolumn{3}{@{}l}{\color{akzent}\bfseries\small IV.~Geometrie und Trigonometrie}\\[1pt]
9 & Elementargeometrie & 77\\
10 & Trigonometrie & 83\\
\multicolumn{3}{@{}l}{\color{akzent}\bfseries\small V.~Die Schwelle zur Analysis}\\[1pt]
11 & Folgen und Grenzwerte & 93\\
12 & Die Ableitung & 99\\
13 & Das Integral & 107\\
\multicolumn{3}{@{}l}{\color{akzent}\bfseries\small VI.~Handwerkszeug}\\[1pt]
14 & Beweisen lernen & 117\\
\multicolumn{3}{@{}l}{\color{akzent}\bfseries\small VII.~Aktueller Stand}\\[1pt]
15 & Der Übergang Schule–Hochschule & 127\\
16 & Die große Formelsammlung & 135\\
17 & Fehlerkatalog & 139\\
18 & Zeittafel & 141\\
19 & Glossar & 143\\
20 & Kommentierte Literatur & 145\\
\end{longtable}}

\section*{Band 1 --- Die Königin der Mathematik}
\addcontentsline{toc}{section}{Band 1 --- Die Königin der Mathematik}
\nopagebreak
{\small\itshape\color{grau} Zahlentheorie $\cdot$ 181 Seiten}\par\vspace{4pt}
\nopagebreak
{\small\renewcommand{\arraystretch}{1.28}
\begin{longtable}{@{}P{0.9cm}P{6.6cm}P{1.0cm}@{}}
\multicolumn{3}{@{}l}{\color{akzent}\bfseries\small I.~Das Fundament}\\[1pt]
1 & Zählen: Die natürlichen Zahlen & 3\\
2 & Teilbarkeit & 11\\
3 & Der größte gemeinsame Teiler & 19\\
4 & Primzahlen und der Fundamentalsatz & 29\\
\multicolumn{3}{@{}l}{\color{akzent}\bfseries\small II.~Die Sprache der Kongruenzen}\\[1pt]
5 & Modulares Rechnen & 43\\
6 & Der chinesische Restsatz & 51\\
7 & Die Sätze von Fermat, Euler und Wilson & 57\\
8 & Quadratische Reste und das Reziprozitätsgesetz & 71\\
\multicolumn{3}{@{}l}{\color{akzent}\bfseries\small III.~Arithmetische Funktionen und die Verteilung der Primzahlen}\\[1pt]
9 & Arithmetische Funktionen & 85\\
10 & Die Verteilung der Primzahlen & 95\\
\multicolumn{3}{@{}l}{\color{akzent}\bfseries\small IV.~Diophantische Gleichungen}\\[1pt]
11 & Pythagoreische Tripel & 107\\
12 & Der unendliche Abstieg und Fermats letzter Satz für n = 4 & 111\\
13 & Kettenbrüche und die Pellsche Gleichung & 115\\
14 & Summen von Quadraten & 121\\
\multicolumn{3}{@{}l}{\color{akzent}\bfseries\small V.~Die Tür zur modernen Zahlentheorie}\\[1pt]
15 & Algebraische Zahlen und Ideale & 127\\
16 & Elliptische Kurven & 131\\
\multicolumn{3}{@{}l}{\color{akzent}\bfseries\small VI.~Anwendungen und aktueller Stand}\\[1pt]
17 & Kryptographie & 141\\
18 & Primzahltests & 145\\
19 & Faktorisierung & 149\\
20 & Der aktuelle Forschungsstand & 153\\
\multicolumn{3}{@{}l}{\color{akzent}\bfseries\small VII.~Anhänge}\\[1pt]
A & Formelsammlung & 161\\
B & Zeitleiste & 165\\
C & Glossar & 167\\
D & Kommentierte Literatur & 169\\
\end{longtable}}

\section*{Band 2 --- Die Kunst der linearen Struktur}
\addcontentsline{toc}{section}{Band 2 --- Die Kunst der linearen Struktur}
\nopagebreak
{\small\itshape\color{grau} Lineare Algebra $\cdot$ 167 Seiten}\par\vspace{4pt}
\nopagebreak
{\small\renewcommand{\arraystretch}{1.28}
\begin{longtable}{@{}P{0.9cm}P{6.6cm}P{1.0cm}@{}}
\multicolumn{3}{@{}l}{\color{akzent}\bfseries\small I.~Das Konkrete: Gleichungen und Matrizen}\\[1pt]
1 & Lineare Gleichungssysteme & 3\\
2 & Matrizen & 15\\
\multicolumn{3}{@{}l}{\color{akzent}\bfseries\small II.~Der abstrakte Sprung}\\[1pt]
3 & Vektorräume & 27\\
4 & Basis und Dimension & 37\\
\multicolumn{3}{@{}l}{\color{akzent}\bfseries\small III.~Lineare Abbildungen}\\[1pt]
5 & Lineare Abbildungen & 49\\
\multicolumn{3}{@{}l}{\color{akzent}\bfseries\small IV.~Determinanten}\\[1pt]
6 & Die Determinante & 63\\
\multicolumn{3}{@{}l}{\color{akzent}\bfseries\small V.~Eigenwerte}\\[1pt]
7 & Eigenwerte und Eigenvektoren & 77\\
8 & Die Jordansche Normalform & 89\\
\multicolumn{3}{@{}l}{\color{akzent}\bfseries\small VI.~Geometrie: Skalarprodukte}\\[1pt]
9 & Skalarprodukte und Orthogonalität & 101\\
10 & Der Spektralsatz & 111\\
11 & Die Singulärwertzerlegung & 119\\
\multicolumn{3}{@{}l}{\color{akzent}\bfseries\small VII.~Numerik und aktueller Stand}\\[1pt]
12 & Numerische lineare Algebra & 129\\
13 & Der aktuelle Forschungsstand & 137\\
\multicolumn{3}{@{}l}{\color{akzent}\bfseries\small VIII.~Anhänge}\\[1pt]
A & Formelsammlung & 147\\
B & Zeitleiste & 151\\
\end{longtable}}

\section*{Band 3 --- Die Kunst des Grenzübergangs}
\addcontentsline{toc}{section}{Band 3 --- Die Kunst des Grenzübergangs}
\nopagebreak
{\small\itshape\color{grau} Analysis $\cdot$ 173 Seiten}\par\vspace{4pt}
\nopagebreak
{\small\renewcommand{\arraystretch}{1.28}
\begin{longtable}{@{}P{0.9cm}P{6.6cm}P{1.0cm}@{}}
\multicolumn{3}{@{}l}{\color{akzent}\bfseries\small I.~Das Fundament: Zahlen, Folgen, Reihen}\\[1pt]
1 & Die reellen Zahlen & 3\\
2 & Folgen und Grenzwerte & 11\\
3 & Reihen & 23\\
\multicolumn{3}{@{}l}{\color{akzent}\bfseries\small II.~Stetigkeit}\\[1pt]
4 & Stetige Funktionen & 37\\
5 & Die elementaren Funktionen & 47\\
\multicolumn{3}{@{}l}{\color{akzent}\bfseries\small III.~Differentialrechnung}\\[1pt]
6 & Die Ableitung & 57\\
7 & Mittelwertsätze und ihre Folgen & 65\\
8 & Der Satz von Taylor & 75\\
\multicolumn{3}{@{}l}{\color{akzent}\bfseries\small IV.~Integralrechnung}\\[1pt]
9 & Das Riemann-Integral & 85\\
10 & Der Hauptsatz und die Integrationstechniken & 93\\
11 & Uneigentliche Integrale & 103\\
\multicolumn{3}{@{}l}{\color{akzent}\bfseries\small V.~Funktionenfolgen}\\[1pt]
12 & Gleichmäßige Konvergenz & 113\\
\multicolumn{3}{@{}l}{\color{akzent}\bfseries\small VI.~Analysis mehrerer Veränderlicher}\\[1pt]
13 & Differentialrechnung im $\R^{n}$ & 125\\
\multicolumn{3}{@{}l}{\color{akzent}\bfseries\small VII.~Ausblick und aktueller Stand}\\[1pt]
14 & Wohin die Analysis führt & 139\\
A & Formelsammlung & 147\\
B & Zeitleiste & 153\\
C & Glossar & 157\\
D & Kommentierte Literatur & 161\\
\end{longtable}}

\section*{Band 4 --- Die Vermessung des Zufalls}
\addcontentsline{toc}{section}{Band 4 --- Die Vermessung des Zufalls}
\nopagebreak
{\small\itshape\color{grau} Wahrscheinlichkeitstheorie $\cdot$ 164 Seiten}\par\vspace{4pt}
\nopagebreak
{\small\renewcommand{\arraystretch}{1.28}
\begin{longtable}{@{}P{0.9cm}P{6.6cm}P{1.0cm}@{}}
\multicolumn{3}{@{}l}{\color{akzent}\bfseries\small I.~Grundlagen}\\[1pt]
1 & Was ist Wahrscheinlichkeit? & 3\\
2 & Kombinatorik und Laplace-Experimente & 13\\
3 & Bedingte Wahrscheinlichkeit und Unabhängigkeit & 21\\
\multicolumn{3}{@{}l}{\color{akzent}\bfseries\small II.~Zufallsvariablen und Verteilungen}\\[1pt]
4 & Zufallsvariablen & 33\\
5 & Die diskreten Verteilungen & 41\\
6 & Die stetigen Verteilungen & 49\\
7 & Erwartungswert, Varianz und Ungleichungen & 57\\
\multicolumn{3}{@{}l}{\color{akzent}\bfseries\small III.~Mehrere Zufallsvariablen}\\[1pt]
8 & Gemeinsame Verteilungen & 69\\
\multicolumn{3}{@{}l}{\color{akzent}\bfseries\small IV.~Die Grenzwertsätze}\\[1pt]
9 & Konvergenzbegriffe & 79\\
10 & Die Gesetze der großen Zahlen & 85\\
11 & Der zentrale Grenzwertsatz & 91\\
\multicolumn{3}{@{}l}{\color{akzent}\bfseries\small V.~Stochastische Prozesse}\\[1pt]
12 & Markow-Ketten & 101\\
13 & Martingale und Irrfahrten & 109\\
\multicolumn{3}{@{}l}{\color{akzent}\bfseries\small VI.~Schließende Statistik}\\[1pt]
14 & Schätzen und Testen & 119\\
\multicolumn{3}{@{}l}{\color{akzent}\bfseries\small VII.~Ausblick und aktueller Stand}\\[1pt]
15 & Der aktuelle Stand & 131\\
A & Formelsammlung & 141\\
B & Verteilungen im Überblick & 145\\
C & Zeitleiste der Wahrscheinlichkeitstheorie & 147\\
D & Glossar & 149\\
E & Kommentierte Literatur & 151\\
\end{longtable}}

\section*{Band 5 --- Die Kunst der Näherung}
\addcontentsline{toc}{section}{Band 5 --- Die Kunst der Näherung}
\nopagebreak
{\small\itshape\color{grau} Numerische Mathematik $\cdot$ 224 Seiten}\par\vspace{4pt}
\nopagebreak
{\small\renewcommand{\arraystretch}{1.28}
\begin{longtable}{@{}P{0.9cm}P{6.6cm}P{1.0cm}@{}}
\multicolumn{3}{@{}l}{\color{akzent}\bfseries\small I.~Rechnen mit endlicher Genauigkeit}\\[1pt]
1 & Zahlen im Rechner & 3\\
2 & Kondition und Stabilität & 17\\
\multicolumn{3}{@{}l}{\color{akzent}\bfseries\small II.~Numerische lineare Algebra}\\[1pt]
3 & Lineare Gleichungssysteme und die LR-Zerlegung & 31\\
4 & QR-Zerlegung und Ausgleichsrechnung & 43\\
5 & Iterative Verfahren für große Systeme & 55\\
6 & Eigenwerte und Singulärwerte & 67\\
\multicolumn{3}{@{}l}{\color{akzent}\bfseries\small III.~Nichtlineare Gleichungen}\\[1pt]
7 & Nullstellen und Fixpunkte & 81\\
\multicolumn{3}{@{}l}{\color{akzent}\bfseries\small IV.~Interpolation und Approximation}\\[1pt]
8 & Polynominterpolation & 97\\
9 & Splines und trigonometrische Interpolation & 109\\
10 & Approximation & 119\\
\multicolumn{3}{@{}l}{\color{akzent}\bfseries\small V.~Integration und Differentiation}\\[1pt]
11 & Numerische Quadratur & 129\\
\multicolumn{3}{@{}l}{\color{akzent}\bfseries\small VI.~Differentialgleichungen}\\[1pt]
12 & Anfangswertprobleme & 147\\
13 & Randwertprobleme und partielle Differentialgleichungen & 159\\
\multicolumn{3}{@{}l}{\color{akzent}\bfseries\small VII.~Optimierung und Zufall}\\[1pt]
14 & Numerische Optimierung und Monte-Carlo & 175\\
15 & Numerik heute & 189\\
\multicolumn{3}{@{}l}{\color{akzent}\bfseries\small IX.~Anhänge}\\[1pt]
A & Formelsammlung & 201\\
B & Verfahren im Überblick & 205\\
C & Zeitleiste der Numerischen Mathematik & 207\\
D & Glossar & 209\\
E & Kommentierte Literatur & 211\\
\end{longtable}}

\section*{Band 6 --- Die Vermessung des Raumes}
\addcontentsline{toc}{section}{Band 6 --- Die Vermessung des Raumes}
\nopagebreak
{\small\itshape\color{grau} Geometrie $\cdot$ 184 Seiten}\par\vspace{4pt}
\nopagebreak
{\small\renewcommand{\arraystretch}{1.28}
\begin{longtable}{@{}P{0.9cm}P{6.6cm}P{1.0cm}@{}}
\multicolumn{3}{@{}l}{\color{akzent}\bfseries\small I.~Klassische Geometrie}\\[1pt]
1 & Euklid und die Axiome & 3\\
2 & Das Dreieck & 11\\
3 & Kreis und Kegelschnitte & 23\\
\multicolumn{3}{@{}l}{\color{akzent}\bfseries\small II.~Analytische Geometrie}\\[1pt]
4 & Koordinaten und Vektoren & 35\\
5 & Quadriken und Hauptachsen & 43\\
\multicolumn{3}{@{}l}{\color{akzent}\bfseries\small III.~Transformationen und Symmetrie}\\[1pt]
6 & Isometrien und Symmetriegruppen & 55\\
7 & Affine und projektive Geometrie & 67\\
\multicolumn{3}{@{}l}{\color{akzent}\bfseries\small IV.~Nichteuklidische Geometrie}\\[1pt]
8 & Das Parallelenproblem & 79\\
9 & Hyperbolische Geometrie & 87\\
10 & Sphärische Geometrie & 97\\
\multicolumn{3}{@{}l}{\color{akzent}\bfseries\small V.~Differentialgeometrie}\\[1pt]
11 & Kurven & 109\\
12 & Flächen & 119\\
13 & Geodäten und Gauß–Bonnet & 129\\
\multicolumn{3}{@{}l}{\color{akzent}\bfseries\small VI.~Polyeder, Konvexität, Packungen}\\[1pt]
14 & Diskrete Geometrie & 139\\
\multicolumn{3}{@{}l}{\color{akzent}\bfseries\small VII.~Aktueller Stand}\\[1pt]
15 & Geometrie heute & 149\\
A & Formelsammlung & 159\\
B & Die Geometrien im Überblick & 165\\
C & Zeitleiste der Geometrie & 167\\
D & Glossar & 169\\
E & Kommentierte Literatur & 171\\
\end{longtable}}

\section*{Band 7 --- Die Wege des Denkens}
\addcontentsline{toc}{section}{Band 7 --- Die Wege des Denkens}
\nopagebreak
{\small\itshape\color{grau} Geschichte der Mathematik $\cdot$ 188 Seiten}\par\vspace{4pt}
\nopagebreak
{\small\renewcommand{\arraystretch}{1.28}
\begin{longtable}{@{}P{0.9cm}P{6.6cm}P{1.0cm}@{}}
\multicolumn{3}{@{}l}{\color{akzent}\bfseries\small I.~Vor dem Beweis}\\[1pt]
1 & Zählen, Zahlzeichen, Rechnen & 3\\
2 & China, Indien, Mesoamerika & 15\\
\multicolumn{3}{@{}l}{\color{akzent}\bfseries\small II.~Das griechische Wunder}\\[1pt]
3 & Die Erfindung des Beweises & 27\\
4 & Euklid und Archimedes & 37\\
5 & Späte Antike und Niedergang & 47\\
\multicolumn{3}{@{}l}{\color{akzent}\bfseries\small III.~Das Erbe wandert}\\[1pt]
6 & Die Mathematik der islamischen Welt & 57\\
7 & Mittelalter und Renaissance in Europa & 67\\
\multicolumn{3}{@{}l}{\color{akzent}\bfseries\small IV.~Die Revolution des 17. Jahrhunderts}\\[1pt]
8 & Koordinaten und Zufall & 81\\
9 & Die Infinitesimalrechnung & 89\\
\multicolumn{3}{@{}l}{\color{akzent}\bfseries\small V.~Das 19. Jahrhundert}\\[1pt]
10 & Die Strenge kommt & 105\\
11 & Die Algebra wird abstrakt & 115\\
12 & Geometrie und Grundlagen & 125\\
\multicolumn{3}{@{}l}{\color{akzent}\bfseries\small VI.~Das 20. Jahrhundert}\\[1pt]
13 & Die Grundlagenkrise & 133\\
14 & Institutionen, Vertreibung, Krieg & 145\\
\multicolumn{3}{@{}l}{\color{akzent}\bfseries\small VII.~Aktueller Stand}\\[1pt]
15 & Mathematikgeschichte heute & 155\\
A & Große Zeittafel & 167\\
B & Kurzbiographien & 171\\
C & Glossar historischer Begriffe & 173\\
D & Kommentierte Literatur & 175\\
\end{longtable}}

\section*{Band 8 --- Die Kunst der Verformung}
\addcontentsline{toc}{section}{Band 8 --- Die Kunst der Verformung}
\nopagebreak
{\small\itshape\color{grau} Topologie $\cdot$ 178 Seiten}\par\vspace{4pt}
\nopagebreak
{\small\renewcommand{\arraystretch}{1.28}
\begin{longtable}{@{}P{0.9cm}P{6.6cm}P{1.0cm}@{}}
\multicolumn{3}{@{}l}{\color{akzent}\bfseries\small I.~Die Grundidee}\\[1pt]
1 & Wo die Topologie herkommt & 3\\
2 & Metrische Räume & 13\\
\multicolumn{3}{@{}l}{\color{akzent}\bfseries\small II.~Topologische Räume}\\[1pt]
3 & Der Raumbegriff & 27\\
4 & Stetigkeit und Konstruktionen & 37\\
\multicolumn{3}{@{}l}{\color{akzent}\bfseries\small III.~Die großen Eigenschaften}\\[1pt]
5 & Zusammenhang & 49\\
6 & Kompaktheit & 55\\
7 & Trennungsaxiome & 65\\
\multicolumn{3}{@{}l}{\color{akzent}\bfseries\small IV.~Algebraische Topologie}\\[1pt]
8 & Homotopie und Fundamentalgruppe & 75\\
9 & Die Ernte: vier klassische Sätze & 85\\
\multicolumn{3}{@{}l}{\color{akzent}\bfseries\small V.~Flächen und Homologie}\\[1pt]
10 & Die Klassifikation der Flächen & 97\\
11 & Homologie & 105\\
\multicolumn{3}{@{}l}{\color{akzent}\bfseries\small VI.~Geometrische und angewandte Topologie}\\[1pt]
12 & Knotentheorie & 117\\
13 & Mannigfaltigkeiten & 125\\
14 & Angewandte Topologie & 133\\
\multicolumn{3}{@{}l}{\color{akzent}\bfseries\small VII.~Aktueller Stand}\\[1pt]
15 & Topologie heute & 145\\
A & Formelsammlung und Übersichten & 153\\
B & Räume und Invarianten im Überblick & 159\\
C & Zeittafel der Topologie & 161\\
D & Glossar & 163\\
E & Kommentierte Literatur & 165\\
\end{longtable}}

\section*{Band 9 --- Die Sprache der Felder}
\addcontentsline{toc}{section}{Band 9 --- Die Sprache der Felder}
\nopagebreak
{\small\itshape\color{grau} Vektor- und Tensoranalysis $\cdot$ 130 Seiten}\par\vspace{4pt}
\nopagebreak
{\small\renewcommand{\arraystretch}{1.28}
\begin{longtable}{@{}P{0.9cm}P{6.6cm}P{1.0cm}@{}}
\multicolumn{3}{@{}l}{\color{akzent}\bfseries\small I.~Felder und Ableitungen}\\[1pt]
1 & Was ein Feld ist & 7\\
2 & Die drei Operatoren & 11\\
3 & Rechenregeln und Identitäten & 19\\
\multicolumn{3}{@{}l}{\color{akzent}\bfseries\small II.~Die Integralsätze}\\[1pt]
4 & Kurven- und Flächenintegrale & 27\\
5 & Green und Gauß & 33\\
6 & Der Satz von Stokes & 37\\
7 & Potentiale und Topologie & 41\\
\multicolumn{3}{@{}l}{\color{akzent}\bfseries\small III.~Krummlinige Koordinaten}\\[1pt]
8 & Andere Koordinaten & 49\\
\multicolumn{3}{@{}l}{\color{akzent}\bfseries\small IV.~Tensoren}\\[1pt]
9 & Was ein Tensor ist & 57\\
10 & Tensoralgebra & 63\\
\multicolumn{3}{@{}l}{\color{akzent}\bfseries\small V.~Tensoranalysis}\\[1pt]
11 & Die kovariante Ableitung & 73\\
12 & Krümmung & 79\\
\multicolumn{3}{@{}l}{\color{akzent}\bfseries\small VI.~Differentialformen}\\[1pt]
13 & Der eine Satz von Stokes & 87\\
14 & Anwendungen & 95\\
\multicolumn{3}{@{}l}{\color{akzent}\bfseries\small VII.~Aktueller Stand}\\[1pt]
15 & Vektor- und Tensoranalysis heute & 103\\
16 & Große Formelsammlung & 111\\
17 & Konventionen und Fallstricke & 115\\
18 & Zeittafel & 117\\
19 & Glossar & 119\\
20 & Kommentierte Literatur & 121\\
\end{longtable}}

\section*{Band 10 --- Die Ordnung des Unendlichen}
\addcontentsline{toc}{section}{Band 10 --- Die Ordnung des Unendlichen}
\nopagebreak
{\small\itshape\color{grau} Logik und Mengenlehre $\cdot$ 124 Seiten}\par\vspace{4pt}
\nopagebreak
{\small\renewcommand{\arraystretch}{1.28}
\begin{longtable}{@{}P{0.9cm}P{6.6cm}P{1.0cm}@{}}
\multicolumn{3}{@{}l}{\color{akzent}\bfseries\small I.~Logik}\\[1pt]
1 & Aussagenlogik & 7\\
2 & Beweiskalküle & 13\\
3 & Prädikatenlogik & 19\\
4 & Vollständigkeit und ihre Folgen & 25\\
\multicolumn{3}{@{}l}{\color{akzent}\bfseries\small II.~Mengenlehre}\\[1pt]
5 & Die naive Mengenlehre und ihr Ende & 33\\
6 & Die Axiome von ZFC & 39\\
7 & Ordinalzahlen & 45\\
8 & Kardinalzahlen & 51\\
\multicolumn{3}{@{}l}{\color{akzent}\bfseries\small III.~Das Auswahlaxiom}\\[1pt]
9 & Wählen ohne Vorschrift & 57\\
\multicolumn{3}{@{}l}{\color{akzent}\bfseries\small IV.~Unentscheidbarkeit}\\[1pt]
10 & Berechenbarkeit & 63\\
11 & Die Unvollständigkeitssätze & 69\\
12 & Weitere Unentscheidbarkeit & 77\\
\multicolumn{3}{@{}l}{\color{akzent}\bfseries\small V.~Modelltheorie und Forcing}\\[1pt]
13 & Modelltheorie & 85\\
14 & Forcing und Unabhängigkeit & 91\\
\multicolumn{3}{@{}l}{\color{akzent}\bfseries\small VI.~Aktueller Stand}\\[1pt]
15 & Logik und Mengenlehre heute & 99\\
\multicolumn{3}{@{}l}{\color{akzent}\bfseries\small VII.~Anhänge}\\[1pt]
16 & Die Axiome im Überblick & 109\\
17 & Die großen Sätze & 111\\
18 & Zeittafel & 113\\
19 & Glossar & 115\\
20 & Kommentierte Literatur & 117\\
\end{longtable}}

\section*{Band 11 --- Die Gestalt der Symmetrie}
\addcontentsline{toc}{section}{Band 11 --- Die Gestalt der Symmetrie}
\nopagebreak
{\small\itshape\color{grau} Algebra $\cdot$ 150 Seiten}\par\vspace{4pt}
\nopagebreak
{\small\renewcommand{\arraystretch}{1.28}
\begin{longtable}{@{}P{0.9cm}P{6.6cm}P{1.0cm}@{}}
\multicolumn{3}{@{}l}{\color{akzent}\bfseries\small I.~Gruppen}\\[1pt]
1 & Was eine Gruppe ist & 7\\
2 & Zyklische Gruppen und Permutationen & 15\\
3 & Nebenklassen und der Satz von Lagrange & 23\\
4 & Normalteiler und Faktorgruppen & 29\\
\multicolumn{3}{@{}l}{\color{akzent}\bfseries\small II.~Die Struktur der Gruppen}\\[1pt]
5 & Gruppenoperationen & 39\\
6 & Die Sylow-Sätze & 47\\
7 & Abelsche Gruppen & 53\\
\multicolumn{3}{@{}l}{\color{akzent}\bfseries\small III.~Ringe}\\[1pt]
8 & Ringe und Ideale & 61\\
9 & Polynomringe & 69\\
\multicolumn{3}{@{}l}{\color{akzent}\bfseries\small IV.~Körper}\\[1pt]
10 & Körpererweiterungen & 79\\
11 & Zerfällungskörper und endliche Körper & 87\\
\multicolumn{3}{@{}l}{\color{akzent}\bfseries\small V.~Galoistheorie}\\[1pt]
12 & Der Hauptsatz & 97\\
13 & Auflösbarkeit durch Radikale & 105\\
\multicolumn{3}{@{}l}{\color{akzent}\bfseries\small VI.~Weiterführende Strukturen}\\[1pt]
14 & Moduln, Darstellungen, Kategorien & 115\\
\multicolumn{3}{@{}l}{\color{akzent}\bfseries\small VII.~Aktueller Stand}\\[1pt]
15 & Algebra heute & 123\\
A & Die Strukturen im Überblick & 131\\
B & Gruppen kleiner Ordnung & 135\\
C & Zeittafel der Algebra & 137\\
D & Glossar & 139\\
E & Kommentierte Literatur & 141\\
\end{longtable}}

\section*{Band 12 --- Die Vollendung der Analysis}
\addcontentsline{toc}{section}{Band 12 --- Die Vollendung der Analysis}
\nopagebreak
{\small\itshape\color{grau} Funktionentheorie $\cdot$ 176 Seiten}\par\vspace{4pt}
\nopagebreak
{\small\renewcommand{\arraystretch}{1.28}
\begin{longtable}{@{}P{0.9cm}P{6.6cm}P{1.0cm}@{}}
\multicolumn{3}{@{}l}{\color{akzent}\bfseries\small I.~Die komplexe Ebene}\\[1pt]
1 & Die komplexen Zahlen & 9\\
2 & Topologie, Folgen und Funktionen & 19\\
\multicolumn{3}{@{}l}{\color{akzent}\bfseries\small II.~Holomorphie}\\[1pt]
3 & Die komplexe Ableitung & 29\\
4 & Potenzreihen und die elementaren Funktionen & 37\\
\multicolumn{3}{@{}l}{\color{akzent}\bfseries\small III.~Integration}\\[1pt]
5 & Kurvenintegrale und der Cauchysche Integralsatz & 51\\
6 & Die Cauchysche Integralformel & 61\\
\multicolumn{3}{@{}l}{\color{akzent}\bfseries\small IV.~Die lokale Struktur}\\[1pt]
7 & Identitätssatz und analytische Fortsetzung & 71\\
8 & Laurentreihen und Singularitäten & 79\\
9 & Der Residuensatz & 87\\
10 & Reelle Integrale mit Residuen & 95\\
\multicolumn{3}{@{}l}{\color{akzent}\bfseries\small V.~Globale Sätze}\\[1pt]
11 & Argumentprinzip und Abbildungsverhalten & 107\\
12 & Konforme Abbildungen & 115\\
\multicolumn{3}{@{}l}{\color{akzent}\bfseries\small VI.~Ganze und meromorphe Funktionen}\\[1pt]
13 & Produktsätze und die Gammafunktion & 125\\
14 & Die Zetafunktion und der Primzahlsatz & 135\\
\multicolumn{3}{@{}l}{\color{akzent}\bfseries\small VII.~Aktueller Stand}\\[1pt]
15 & Funktionentheorie heute & 147\\
A & Formelsammlung & 155\\
B & Standardintegrale und Wege & 159\\
C & Die großen Sätze im Überblick & 161\\
D & Zeittafel der Funktionentheorie & 163\\
E & Glossar & 165\\
\end{longtable}}

\section*{Band 13 --- Das Maß und der Raum}
\addcontentsline{toc}{section}{Band 13 --- Das Maß und der Raum}
\nopagebreak
{\small\itshape\color{grau} Maßtheorie und Funktionalanalysis $\cdot$ 175 Seiten}\par\vspace{4pt}
\nopagebreak
{\small\renewcommand{\arraystretch}{1.28}
\begin{longtable}{@{}P{0.9cm}P{6.6cm}P{1.0cm}@{}}
\multicolumn{3}{@{}l}{\color{akzent}\bfseries\small I.~Warum Riemann nicht genügt}\\[1pt]
1 & Die drei Defekte des Riemann-Integrals & 9\\
2 & Das Problem des Messens & 19\\
\multicolumn{3}{@{}l}{\color{akzent}\bfseries\small II.~Maß und Integral}\\[1pt]
3 & $\sigma$-Algebren und Maße & 29\\
4 & Das Lebesgue-Maß & 37\\
5 & Messbare Funktionen und das Integral & 45\\
6 & Die drei Konvergenzsätze & 53\\
7 & Produktmaße, Fubini und Transformation & 63\\
\multicolumn{3}{@{}l}{\color{akzent}\bfseries\small III.~Das Scharnier: die $L^{p}$ -Räume}\\[1pt]
8 & Die Räume $L^{p}$ & 73\\
9 & Radon-Nikodym und Dualität & 83\\
\multicolumn{3}{@{}l}{\color{akzent}\bfseries\small IV.~Banachräume}\\[1pt]
10 & Normierte Räume und Hahn-Banach & 93\\
11 & Baire und die drei Prinzipien & 101\\
\multicolumn{3}{@{}l}{\color{akzent}\bfseries\small V.~Hilberträume}\\[1pt]
12 & Geometrie im Unendlichdimensionalen & 111\\
\multicolumn{3}{@{}l}{\color{akzent}\bfseries\small VI.~Operatoren}\\[1pt]
13 & Kompakte Operatoren und Fredholm & 121\\
14 & Der Spektralsatz und die Quantenmechanik & 129\\
\multicolumn{3}{@{}l}{\color{akzent}\bfseries\small VII.~Aktueller Stand}\\[1pt]
15 & Wo die Forschung heute steht & 147\\
A & Formelsammlung und Übersicht & 157\\
B & Kommentierte Literatur & 167\\
\end{longtable}}

\section*{Band 14 --- Die Gesetze der Veränderung}
\addcontentsline{toc}{section}{Band 14 --- Die Gesetze der Veränderung}
\nopagebreak
{\small\itshape\color{grau} Differentialgleichungen $\cdot$ 214 Seiten}\par\vspace{4pt}
\nopagebreak
{\small\renewcommand{\arraystretch}{1.28}
\begin{longtable}{@{}P{0.9cm}P{6.6cm}P{1.0cm}@{}}
\multicolumn{3}{@{}l}{\color{akzent}\bfseries\small I.~Was eine Differentialgleichung ist}\\[1pt]
1 & Der Begriff und das Richtungsfeld & 9\\
2 & Die elementaren Lösungsverfahren & 19\\
\multicolumn{3}{@{}l}{\color{akzent}\bfseries\small II.~Die Theorie erster Ordnung}\\[1pt]
3 & Existenz und Eindeutigkeit & 33\\
4 & Lineare Systeme und das Matrixexponential & 45\\
\multicolumn{3}{@{}l}{\color{akzent}\bfseries\small III.~Höhere Ordnung und Schwingungen}\\[1pt]
5 & Lineare Gleichungen mit konstanten Koeffizienten & 57\\
6 & Potenzreihen und spezielle Funktionen & 67\\
7 & Randwertprobleme und Sturm-Liouville & 79\\
\multicolumn{3}{@{}l}{\color{akzent}\bfseries\small IV.~Qualitative Theorie}\\[1pt]
8 & Phasenporträt und Stabilität & 93\\
9 & Grenzzyklen, Bifurkationen und Chaos & 105\\
\multicolumn{3}{@{}l}{\color{akzent}\bfseries\small V.~Partielle Differentialgleichungen}\\[1pt]
10 & Klassifikation und Charakteristiken & 121\\
11 & Die Wellengleichung & 131\\
12 & Die Wärmeleitungsgleichung & 141\\
13 & Die Laplacegleichung & 153\\
\multicolumn{3}{@{}l}{\color{akzent}\bfseries\small VI.~Moderne Methoden}\\[1pt]
14 & Schwache Lösungen und numerische Verfahren & 167\\
\multicolumn{3}{@{}l}{\color{akzent}\bfseries\small VII.~Aktueller Stand}\\[1pt]
15 & Wo die Forschung heute steht & 183\\
A & Formelsammlung und Übersicht & 195\\
B & Kommentierte Literatur & 205\\
\end{longtable}}

\section*{Band 15 --- Die Kunst des Zählens}
\addcontentsline{toc}{section}{Band 15 --- Die Kunst des Zählens}
\nopagebreak
{\small\itshape\color{grau} Kombinatorik $\cdot$ 200 Seiten}\par\vspace{4pt}
\nopagebreak
{\small\renewcommand{\arraystretch}{1.28}
\begin{longtable}{@{}P{0.9cm}P{6.6cm}P{1.0cm}@{}}
\multicolumn{3}{@{}l}{\color{akzent}\bfseries\small I.~Grundlegendes Zählen}\\[1pt]
1 & Die Grundregeln & 9\\
2 & Die vier Grundaufgaben & 19\\
\multicolumn{3}{@{}l}{\color{akzent}\bfseries\small II.~Verfeinerungen}\\[1pt]
3 & Das Pascalsche Dreieck & 33\\
4 & Inklusion und Exklusion & 43\\
\multicolumn{3}{@{}l}{\color{akzent}\bfseries\small III.~Erzeugende Funktionen}\\[1pt]
5 & Gewöhnliche erzeugende Funktionen & 55\\
6 & Exponentielle erzeugende Funktionen & 67\\
7 & Rekursionen, Catalanzahlen, Asymptotik & 77\\
\multicolumn{3}{@{}l}{\color{akzent}\bfseries\small IV.~Struktur des Zählens}\\[1pt]
8 & Der Zwölffachweg & 89\\
9 & Symmetrie und Zählen & 99\\
\multicolumn{3}{@{}l}{\color{akzent}\bfseries\small V.~Graphentheorie}\\[1pt]
10 & Graphen und Bäume & 111\\
11 & Planarität und Färbung & 121\\
12 & Matchings, Flüsse und der Satz von Hall & 133\\
\multicolumn{3}{@{}l}{\color{akzent}\bfseries\small VI.~Ramseytheorie und die probabilistische Methode}\\[1pt]
13 & Ramseytheorie & 145\\
14 & Die probabilistische Methode und extremale Kombinatorik & 155\\
\multicolumn{3}{@{}l}{\color{akzent}\bfseries\small VII.~Aktueller Stand}\\[1pt]
15 & Wo die Forschung heute steht & 169\\
A & Formelsammlung und Übersicht & 179\\
B & Kommentierte Literatur & 191\\
\end{longtable}}

\section*{Band 16 --- Punkte und Linien}
\addcontentsline{toc}{section}{Band 16 --- Punkte und Linien}
\nopagebreak
{\small\itshape\color{grau} Graphentheorie $\cdot$ 270 Seiten}\par\vspace{4pt}
\nopagebreak
{\small\renewcommand{\arraystretch}{1.28}
\begin{longtable}{@{}P{0.9cm}P{6.6cm}P{1.0cm}@{}}
\multicolumn{3}{@{}l}{\color{akzent}\bfseries\small I.~Der Begriff}\\[1pt]
1 & Was ein Graph ist & 9\\
2 & Wege, Kreise und Zusammenhang & 33\\
3 & Wie man Graphen aufschreibt & 53\\
\multicolumn{3}{@{}l}{\color{akzent}\bfseries\small II.~Bäume}\\[1pt]
4 & Bäume und Wälder & 67\\
5 & Spannbäume und wie man sie zählt & 81\\
\multicolumn{3}{@{}l}{\color{akzent}\bfseries\small III.~Durchlaufen}\\[1pt]
6 & Eulerzüge & 105\\
7 & Hamiltonkreise & 117\\
\multicolumn{3}{@{}l}{\color{akzent}\bfseries\small IV.~Zeichnen und Färben}\\[1pt]
8 & Planare Graphen & 137\\
9 & Färbungen & 151\\
10 & Das chromatische Polynom & 165\\
\multicolumn{3}{@{}l}{\color{akzent}\bfseries\small V.~Paarungen und Flüsse}\\[1pt]
11 & Paarungen & 179\\
12 & Flüsse und Schnitte & 193\\
\multicolumn{3}{@{}l}{\color{akzent}\bfseries\small VI.~Grenzen und Zufall}\\[1pt]
13 & Extremale Graphentheorie & 209\\
14 & Spektren und Zufallsgraphen & 223\\
\multicolumn{3}{@{}l}{\color{akzent}\bfseries\small VII.~Aktueller Stand}\\[1pt]
15 & Wo die Forschung heute steht & 239\\
\multicolumn{3}{@{}l}{\color{akzent}\bfseries\small VIII.~Anhang}\\[1pt]
A & Formelsammlung und Übersicht & 251\\
B & Kommentierte Quellen & 259\\
\end{longtable}}

\section*{Band 17 --- Was die Zahlen hergeben}
\addcontentsline{toc}{section}{Band 17 --- Was die Zahlen hergeben}
\nopagebreak
{\small\itshape\color{grau} Statistik $\cdot$ 282 Seiten}\par\vspace{4pt}
\nopagebreak
{\small\renewcommand{\arraystretch}{1.28}
\begin{longtable}{@{}P{0.9cm}P{6.6cm}P{1.0cm}@{}}
\multicolumn{3}{@{}l}{\color{akzent}\bfseries\small I.~Beschreiben}\\[1pt]
1 & Daten, Skalen und die erste Falle & 9\\
2 & Lage und Streuung & 25\\
3 & Zusammenhang & 45\\
\multicolumn{3}{@{}l}{\color{akzent}\bfseries\small II.~Schätzen}\\[1pt]
4 & Von der Stichprobe zur Aussage & 63\\
5 & Punktschätzer und ihre Güte & 79\\
6 & Konfidenzintervalle & 95\\
\multicolumn{3}{@{}l}{\color{akzent}\bfseries\small III.~Testen}\\[1pt]
7 & Die Logik des Tests & 113\\
8 & Die klassischen Tests & 131\\
9 & Verteilungsfreie Verfahren und Simulation & 149\\
\multicolumn{3}{@{}l}{\color{akzent}\bfseries\small IV.~Modelle}\\[1pt]
10 & Lineare Regression & 167\\
11 & Andere Zielgrößen & 185\\
\multicolumn{3}{@{}l}{\color{akzent}\bfseries\small V.~Medizinische Statistik}\\[1pt]
12 & Studiendesign & 203\\
13 & Diagnostische Tests & 221\\
14 & Viele Tests auf einmal & 235\\
\multicolumn{3}{@{}l}{\color{akzent}\bfseries\small VI.~Aktueller Stand}\\[1pt]
15 & Wo die Statistik heute steht & 253\\
\multicolumn{3}{@{}l}{\color{akzent}\bfseries\small VII.~Anhang}\\[1pt]
A & Formelsammlung und Übersicht & 263\\
B & Kommentierte Quellen & 271\\
\end{longtable}}

\section*{Band 18 --- Wachsen und Vergehen}
\addcontentsline{toc}{section}{Band 18 --- Wachsen und Vergehen}
\nopagebreak
{\small\itshape\color{grau} Biomathematik $\cdot$ 262 Seiten}\par\vspace{4pt}
\nopagebreak
{\small\renewcommand{\arraystretch}{1.28}
\begin{longtable}{@{}P{0.9cm}P{6.6cm}P{1.0cm}@{}}
\multicolumn{3}{@{}l}{\color{akzent}\bfseries\small I.~Wachstum}\\[1pt]
1 & Exponentielles Wachstum & 9\\
2 & Beschränktes Wachstum & 23\\
3 & Diskrete Modelle und Chaos & 37\\
\multicolumn{3}{@{}l}{\color{akzent}\bfseries\small II.~Wechselwirkung}\\[1pt]
4 & Räuber und Beute & 53\\
5 & Wettbewerb und Koexistenz & 69\\
6 & Stabilität & 83\\
\multicolumn{3}{@{}l}{\color{akzent}\bfseries\small III.~Ausbreitung}\\[1pt]
7 & Das SIR-Modell & 101\\
8 & Impfung und Schwellenwerte & 117\\
9 & Netzwerke und Heterogenität & 131\\
\multicolumn{3}{@{}l}{\color{akzent}\bfseries\small IV.~Der Körper}\\[1pt]
10 & Pharmakokinetik & 147\\
11 & Enzym- und Rezeptorkinetik & 165\\
12 & Erregbare Zellen & 181\\
\multicolumn{3}{@{}l}{\color{akzent}\bfseries\small V.~Krankheit}\\[1pt]
13 & Tumorwachstum und Therapie & 199\\
14 & Resistenz und Evolution & 213\\
\multicolumn{3}{@{}l}{\color{akzent}\bfseries\small VI.~Aktueller Stand}\\[1pt]
15 & Was Modelle leisten & 231\\
\multicolumn{3}{@{}l}{\color{akzent}\bfseries\small VII.~Anhang}\\[1pt]
A & Formelsammlung und Übersicht & 241\\
B & Kommentierte Quellen & 251\\
\end{longtable}}

\section*{Band 19 --- Das Beste finden}
\addcontentsline{toc}{section}{Band 19 --- Das Beste finden}
\nopagebreak
{\small\itshape\color{grau} Optimierung $\cdot$ 252 Seiten}\par\vspace{4pt}
\nopagebreak
{\small\renewcommand{\arraystretch}{1.28}
\begin{longtable}{@{}P{0.9cm}P{6.6cm}P{1.0cm}@{}}
\multicolumn{3}{@{}l}{\color{akzent}\bfseries\small I.~Lineare Probleme}\\[1pt]
1 & Was heißt optimal? & 7\\
2 & Lineare Probleme & 21\\
3 & Das Simplexverfahren & 37\\
\multicolumn{3}{@{}l}{\color{akzent}\bfseries\small II.~Dualität}\\[1pt]
4 & Das duale Problem & 53\\
5 & Schattenpreise und Sensitivität & 69\\
6 & Netzwerke und Ganzzahligkeit & 83\\
\multicolumn{3}{@{}l}{\color{akzent}\bfseries\small III.~Ohne Linearität}\\[1pt]
7 & Konvexität & 103\\
8 & Lagrange und die KKT-Bedingungen & 117\\
9 & Abstiegsverfahren & 133\\
\multicolumn{3}{@{}l}{\color{akzent}\bfseries\small IV.~Ganzzahligkeit}\\[1pt]
10 & Wenn Bruchteile nicht gehen & 151\\
11 & Verzweigen und Beschränken & 167\\
12 & Was schwer ist und warum & 181\\
\multicolumn{3}{@{}l}{\color{akzent}\bfseries\small V.~Unsicherheit}\\[1pt]
13 & Wenn die Daten wackeln & 197\\
14 & Mehrere Ziele & 209\\
\multicolumn{3}{@{}l}{\color{akzent}\bfseries\small VI.~Aktueller Stand}\\[1pt]
15 & Was Optimierung leistet & 223\\
\multicolumn{3}{@{}l}{\color{akzent}\bfseries\small VII.~Anhang}\\[1pt]
A & Formelsammlung und Übersicht & 231\\
B & Kommentierte Quellen & 241\\
\end{longtable}}

\section*{Band 20 --- Was Nachrichten wiegen}
\addcontentsline{toc}{section}{Band 20 --- Was Nachrichten wiegen}
\nopagebreak
{\small\itshape\color{grau} Information und Kryptographie $\cdot$ 242 Seiten}\par\vspace{4pt}
\nopagebreak
{\small\renewcommand{\arraystretch}{1.28}
\begin{longtable}{@{}P{0.9cm}P{6.6cm}P{1.0cm}@{}}
\multicolumn{3}{@{}l}{\color{akzent}\bfseries\small I.~Information messen}\\[1pt]
1 & Was Information ist & 9\\
2 & Entropie & 23\\
3 & Quellencodierung & 37\\
\multicolumn{3}{@{}l}{\color{akzent}\bfseries\small II.~Der gestörte Kanal}\\[1pt]
4 & Transinformation und Kapazität & 55\\
5 & Das Kanalcodierungstheorem & 69\\
6 & Fehlerkorrigierende Codes & 83\\
\multicolumn{3}{@{}l}{\color{akzent}\bfseries\small III.~Zufall und Komplexität}\\[1pt]
7 & Kolmogorow-Komplexität & 99\\
8 & Information in der Statistik & 113\\
\multicolumn{3}{@{}l}{\color{akzent}\bfseries\small IV.~Geheimhaltung}\\[1pt]
9 & Perfekte Sicherheit & 129\\
10 & Blockchiffren & 141\\
11 & Hashfunktionen und Unversehrtheit & 153\\
\multicolumn{3}{@{}l}{\color{akzent}\bfseries\small V.~Öffentliche Schlüssel}\\[1pt]
12 & Schlüsseltausch über offene Leitung & 169\\
13 & RSA & 181\\
14 & Elliptische Kurven und Signaturen & 195\\
\multicolumn{3}{@{}l}{\color{akzent}\bfseries\small VI.~Aktueller Stand}\\[1pt]
15 & Was Information leistet & 211\\
\multicolumn{3}{@{}l}{\color{akzent}\bfseries\small VII.~Anhang}\\[1pt]
A & Formelsammlung und Übersicht & 221\\
B & Kommentierte Quellen & 231\\
\end{longtable}}

\section*{Band 21 --- Wenn andere mitrechnen}
\addcontentsline{toc}{section}{Band 21 --- Wenn andere mitrechnen}
\nopagebreak
{\small\itshape\color{grau} Spieltheorie $\cdot$ 236 Seiten}\par\vspace{4pt}
\nopagebreak
{\small\renewcommand{\arraystretch}{1.28}
\begin{longtable}{@{}P{0.9cm}P{6.6cm}P{1.0cm}@{}}
\multicolumn{3}{@{}l}{\color{akzent}\bfseries\small I.~Entgegengesetzte Interessen}\\[1pt]
1 & Was ein Spiel ist & 7\\
2 & Der Minimax-Satz & 21\\
3 & Spiele und lineare Optimierung & 35\\
\multicolumn{3}{@{}l}{\color{akzent}\bfseries\small II.~Gleichgewichte}\\[1pt]
4 & Das Nash-Gleichgewicht & 51\\
5 & Drei Spiele, die alles erklären & 63\\
6 & Existenz und Berechnung & 75\\
\multicolumn{3}{@{}l}{\color{akzent}\bfseries\small III.~Über die Zeit}\\[1pt]
7 & Reihenfolge und Glaubwürdigkeit & 91\\
8 & Wiederholte Spiele & 103\\
\multicolumn{3}{@{}l}{\color{akzent}\bfseries\small IV.~Evolution}\\[1pt]
9 & Evolutionär stabile Strategien & 119\\
10 & Replikatordynamik & 131\\
\multicolumn{3}{@{}l}{\color{akzent}\bfseries\small V.~Regeln entwerfen}\\[1pt]
11 & Kollektivgüter und soziale Dilemmata & 147\\
12 & Abstimmen & 161\\
13 & Auktionen und Zuteilung & 175\\
14 & Verhandlung und Koalitionen & 189\\
\multicolumn{3}{@{}l}{\color{akzent}\bfseries\small VI.~Aktueller Stand}\\[1pt]
15 & Was die Theorie leistet & 205\\
\multicolumn{3}{@{}l}{\color{akzent}\bfseries\small VII.~Anhang}\\[1pt]
A & Formelsammlung und Übersicht & 215\\
B & Kommentierte Quellen & 225\\
\multicolumn{3}{@{}l}{\color{akzent}\bfseries\small VI.~Stand}\\[1pt]
\end{longtable}}

\section*{Band 22 --- Was sich verallgemeinern lässt}
\addcontentsline{toc}{section}{Band 22 --- Was sich verallgemeinern lässt}
\nopagebreak
{\small\itshape\color{grau} Maschinelles Lernen $\cdot$ 204 Seiten}\par\vspace{4pt}
\nopagebreak
{\small\renewcommand{\arraystretch}{1.28}
\begin{longtable}{@{}P{0.9cm}P{6.6cm}P{1.0cm}@{}}
\multicolumn{3}{@{}l}{\color{akzent}\bfseries\small I.~Das Problem}\\[1pt]
1 & Was Lernen heißt & 7\\
2 & Zu wenig und zu viel & 21\\
3 & Wie man prüft --- und wie man sich betrügt & 33\\
\multicolumn{3}{@{}l}{\color{akzent}\bfseries\small II.~Lineare Verfahren}\\[1pt]
4 & Die Freiheit begrenzen & 47\\
5 & Eine Wahrscheinlichkeit vorhersagen & 57\\
6 & Der breiteste Streifen & 69\\
\multicolumn{3}{@{}l}{\color{akzent}\bfseries\small III.~Die Theorie}\\[1pt]
7 & Wie viele Beispiele man braucht & 85\\
8 & Warum Einfaches verallgemeinert & 97\\
\multicolumn{3}{@{}l}{\color{akzent}\bfseries\small IV.~Nichtlineare Verfahren}\\[1pt]
9 & Bäume und Wälder & 109\\
10 & Ähnliche Fälle und der Fluch der Dimension & 121\\
11 & Neuronale Netze & 131\\
\multicolumn{3}{@{}l}{\color{akzent}\bfseries\small V.~Ohne Etiketten, in der Tiefe}\\[1pt]
12 & Gruppen finden & 145\\
13 & Die dünne Fläche & 155\\
14 & Das Rätsel der tiefen Netze & 165\\
\multicolumn{3}{@{}l}{\color{akzent}\bfseries\small VI.~Aktueller Stand}\\[1pt]
15 & Was sich verallgemeinern lässt & 177\\
\multicolumn{3}{@{}l}{\color{akzent}\bfseries\small VII.~Anhang}\\[1pt]
A & Formelsammlung und Übersicht & 185\\
B & Kommentierte Quellen & 193\\
\end{longtable}}

\section*{Band 23 --- Woran es liegt}
\addcontentsline{toc}{section}{Band 23 --- Woran es liegt}
\nopagebreak
{\small\itshape\color{grau} Kausale Schlüsse $\cdot$ 158 Seiten}\par\vspace{4pt}
\nopagebreak
{\small\renewcommand{\arraystretch}{1.28}
\begin{longtable}{@{}P{0.9cm}P{6.6cm}P{1.0cm}@{}}
\multicolumn{3}{@{}l}{\color{akzent}\bfseries\small I.~Das Problem}\\[1pt]
1 & Zusammenhang und Ursache & 7\\
2 & Was gewesen wäre & 17\\
3 & Das Los & 27\\
\multicolumn{3}{@{}l}{\color{akzent}\bfseries\small II.~Graphen und Regeln}\\[1pt]
4 & Annahmen aufschreiben & 39\\
5 & Was man herausrechnen darf & 49\\
6 & Der Hintertürtest & 59\\
\multicolumn{3}{@{}l}{\color{akzent}\bfseries\small III.~Anpassen}\\[1pt]
7 & Eine Zahl statt zwanzig & 69\\
8 & Gewichten und doppelte Absicherung & 77\\
\multicolumn{3}{@{}l}{\color{akzent}\bfseries\small IV.~Wenn Anpassen nicht reicht}\\[1pt]
9 & Instrumente & 87\\
10 & Natürliche Versuche & 95\\
11 & Die Studie nachbilden & 103\\
\multicolumn{3}{@{}l}{\color{akzent}\bfseries\small V.~Über die Zeit, im Einzelfall}\\[1pt]
12 & Behandeln über die Zeit & 113\\
13 & Auf welchem Weg & 121\\
14 & Für wen es wirkt & 127\\
\multicolumn{3}{@{}l}{\color{akzent}\bfseries\small VI.~Aktueller Stand}\\[1pt]
15 & Woran es liegt & 137\\
\multicolumn{3}{@{}l}{\color{akzent}\bfseries\small VII.~Anhang}\\[1pt]
A & Formelsammlung und Übersicht & 143\\
B & Kommentierte Quellen & 149\\
\end{longtable}}

\section*{Band 24 --- Was zu tun ist}
\addcontentsline{toc}{section}{Band 24 --- Was zu tun ist}
\nopagebreak
{\small\itshape\color{grau} Entscheidungen unter Unsicherheit $\cdot$ 140 Seiten}\par\vspace{4pt}
\nopagebreak
{\small\renewcommand{\arraystretch}{1.28}
\begin{longtable}{@{}P{0.9cm}P{6.6cm}P{1.0cm}@{}}
\multicolumn{3}{@{}l}{\color{akzent}\bfseries\small I.~Grundlagen}\\[1pt]
1 & Abwägen & 7\\
2 & Woher der Nutzen kommt & 15\\
3 & Was Wissen wert ist & 23\\
\multicolumn{3}{@{}l}{\color{akzent}\bfseries\small II.~Klinische Entscheidungen}\\[1pt]
4 & Zwei Schwellen & 33\\
5 & Bäume und ihre Empfindlichkeit & 41\\
6 & Über zwanzig Jahre & 47\\
\multicolumn{3}{@{}l}{\color{akzent}\bfseries\small III.~Über die Zeit}\\[1pt]
7 & Rückwärts denken & 57\\
8 & Ohne Ende & 63\\
9 & Erkunden oder ausnutzen & 69\\
\multicolumn{3}{@{}l}{\color{akzent}\bfseries\small IV.~Grenzen des Modells}\\[1pt]
10 & Mehr als ein Ziel & 77\\
11 & Unsichere Wahrscheinlichkeiten & 83\\
12 & Wie Menschen wirklich entscheiden & 91\\
\multicolumn{3}{@{}l}{\color{akzent}\bfseries\small V.~Entscheiden für andere}\\[1pt]
13 & Wenn mehrere entscheiden & 101\\
14 & Entscheiden für viele & 109\\
\multicolumn{3}{@{}l}{\color{akzent}\bfseries\small VI.~Aktueller Stand}\\[1pt]
15 & Was zu tun ist & 117\\
\multicolumn{3}{@{}l}{\color{akzent}\bfseries\small VII.~Anhang}\\[1pt]
A & Formelsammlung und Übersicht & 125\\
B & Kommentierte Quellen & 131\\
\end{longtable}}

\section*{Band 25 --- Wenn der Zufall weitergeht}
\addcontentsline{toc}{section}{Band 25 --- Wenn der Zufall weitergeht}
\nopagebreak
{\small\itshape\color{grau} Stochastische Prozesse $\cdot$ 158 Seiten}\par\vspace{4pt}
\nopagebreak
{\small\renewcommand{\arraystretch}{1.28}
\begin{longtable}{@{}P{0.9cm}P{6.6cm}P{1.0cm}@{}}
\multicolumn{3}{@{}l}{\color{akzent}\bfseries\small I.~Ketten}\\[1pt]
1 & Die Irrfahrt & 7\\
2 & Ketten ohne Gedächtnis & 17\\
3 & Wohin die Kette strebt & 27\\
\multicolumn{3}{@{}l}{\color{akzent}\bfseries\small II.~Sprünge in stetiger Zeit}\\[1pt]
4 & Ereignisse ohne Gedächtnis & 37\\
5 & Warten & 45\\
6 & Aussterben oder Ausbruch & 53\\
\multicolumn{3}{@{}l}{\color{akzent}\bfseries\small III.~Faire Spiele}\\[1pt]
7 & Was man schon weiß & 63\\
8 & Faire Spiele & 71\\
9 & Wohin faire Spiele laufen & 79\\
\multicolumn{3}{@{}l}{\color{akzent}\bfseries\small IV.~Die Brownsche Bewegung}\\[1pt]
10 & Stetig und nirgends glatt & 89\\
11 & Ein Integral für zerknitterte Pfade & 99\\
12 & Die Kettenregel mit Zusatz & 107\\
\multicolumn{3}{@{}l}{\color{akzent}\bfseries\small V.~Anwendungen}\\[1pt]
13 & Was eine Option wert ist & 117\\
14 & Prozesse in Klinik und Natur & 125\\
\multicolumn{3}{@{}l}{\color{akzent}\bfseries\small VI.~Aktueller Stand}\\[1pt]
15 & Was der Zufall mitbringt & 135\\
\multicolumn{3}{@{}l}{\color{akzent}\bfseries\small VII.~Anhang}\\[1pt]
A & Formelsammlung und Übersicht & 143\\
B & Kommentierte Quellen & 149\\
\multicolumn{3}{@{}l}{\color{akzent}\bfseries\small VI.~Stand}\\[1pt]
\end{longtable}}

\section*{Band 26 --- Wenn die Daten sich selbst bestimmen}
\addcontentsline{toc}{section}{Band 26 --- Wenn die Daten sich selbst bestimmen}
\nopagebreak
{\small\itshape\color{grau} Ökonometrie $\cdot$ 156 Seiten}\par\vspace{4pt}
\nopagebreak
{\small\renewcommand{\arraystretch}{1.28}
\begin{longtable}{@{}P{0.9cm}P{6.6cm}P{1.0cm}@{}}
\multicolumn{3}{@{}l}{\color{akzent}\bfseries\small I.~Das Problem}\\[1pt]
1 & Wenn der Regressor nicht unschuldig ist & 7\\
2 & Was die Daten nicht hergeben & 15\\
3 & Eine Veränderung von außen & 23\\
\multicolumn{3}{@{}l}{\color{akzent}\bfseries\small II.~Querschnittsdaten}\\[1pt]
4 & Momente statt Verteilungen & 33\\
5 & Ja oder nein, null oder mehr & 41\\
6 & Wer in den Daten ist & 49\\
\multicolumn{3}{@{}l}{\color{akzent}\bfseries\small III.~Dieselben über die Zeit}\\[1pt]
7 & Jeder sein eigener Vergleich & 59\\
8 & Wenn gestern heute erklärt & 67\\
9 & Vorher, nachher, und mit wem & 75\\
\multicolumn{3}{@{}l}{\color{akzent}\bfseries\small IV.~Zeitreihen}\\[1pt]
10 & Eine Reihe, viele Zeitpunkte & 85\\
11 & Zusammenhänge, die keine sind & 93\\
12 & Zwei Irrfahrten, ein Band & 101\\
\multicolumn{3}{@{}l}{\color{akzent}\bfseries\small V.~Vorhersage und Politik}\\[1pt]
13 & Was ein gutes Modell nicht kann & 113\\
14 & Der Leitfall, zu Ende gerechnet & 121\\
\multicolumn{3}{@{}l}{\color{akzent}\bfseries\small VI.~Aktueller Stand}\\[1pt]
15 & Was die Ökonometrie beiträgt & 131\\
\multicolumn{3}{@{}l}{\color{akzent}\bfseries\small VII.~Anhang}\\[1pt]
A & Formelsammlung und Übersicht & 139\\
B & Kommentierte Quellen & 147\\
\end{longtable}}

\section*{Band 27 --- Was der Rechner vom Zufall weiß}
\addcontentsline{toc}{section}{Band 27 --- Was der Rechner vom Zufall weiß}
\nopagebreak
{\small\itshape\color{grau} Numerische Simulation $\cdot$ 170 Seiten}\par\vspace{4pt}
\nopagebreak
{\small\renewcommand{\arraystretch}{1.28}
\begin{longtable}{@{}P{0.9cm}P{6.6cm}P{1.0cm}@{}}
\multicolumn{3}{@{}l}{\color{akzent}\bfseries\small I.~Der Zufall im Rechner}\\[1pt]
1 & Ein Experiment im Rechner & 7\\
2 & Zufall aus einer Formel & 15\\
3 & Der Fluch der Dimension & 25\\
\multicolumn{3}{@{}l}{\color{akzent}\bfseries\small II.~Genauer werden}\\[1pt]
4 & Den Fehler richtig messen & 35\\
5 & Weniger Streuung ohne mehr Läufe & 45\\
6 & Dort ziehen, wo etwas passiert & 55\\
\multicolumn{3}{@{}l}{\color{akzent}\bfseries\small III.~Ketten und Pfade}\\[1pt]
7 & Ziehen, ohne ziehen zu können & 67\\
8 & Ist die Kette angekommen? & 77\\
9 & Fliegen statt kriechen, und Pfade rechnen & 85\\
\multicolumn{3}{@{}l}{\color{akzent}\bfseries\small IV.~Systeme über die Zeit}\\[1pt]
10 & Die Station als Programm & 97\\
11 & Was hilft, und wie selten es schlimm wird & 107\\
12 & Was jetzt gerade los ist & 115\\
\multicolumn{3}{@{}l}{\color{akzent}\bfseries\small V.~Die Praxis}\\[1pt]
13 & Dasselbe noch einmal & 125\\
14 & Was man dem Ergebnis glauben darf & 135\\
\multicolumn{3}{@{}l}{\color{akzent}\bfseries\small VI.~Aktueller Stand}\\[1pt]
15 & Was der Rechner vom Zufall weiß & 145\\
\multicolumn{3}{@{}l}{\color{akzent}\bfseries\small VII.~Anhang}\\[1pt]
A & Formelsammlung und Übersicht & 153\\
B & Kommentierte Quellen & 161\\
\end{longtable}}

\section*{Band 28 --- Wo die Vorhersage endet}
\addcontentsline{toc}{section}{Band 28 --- Wo die Vorhersage endet}
\nopagebreak
{\small\itshape\color{grau} Dynamische Systeme und Chaos $\cdot$ 166 Seiten}\par\vspace{4pt}
\nopagebreak
{\small\renewcommand{\arraystretch}{1.28}
\begin{longtable}{@{}P{0.9cm}P{6.6cm}P{1.0cm}@{}}
\multicolumn{3}{@{}l}{\color{akzent}\bfseries\small I.~Der Fluss}\\[1pt]
1 & Ein Zustand und seine Zukunft & 7\\
2 & Ruhe, und ob sie hält & 15\\
3 & Die Schwingung, die von allein kommt & 25\\
\multicolumn{3}{@{}l}{\color{akzent}\bfseries\small II.~Wenn Parameter wandern}\\[1pt]
4 & Wenn aus einem Fixpunkt drei werden & 37\\
5 & Wie aus der Ruhe ein Takt wird & 47\\
6 & Ein Takt von außen & 55\\
\multicolumn{3}{@{}l}{\color{akzent}\bfseries\small III.~Der Weg ins Chaos}\\[1pt]
7 & Wenn die Periode sich verdoppelt & 67\\
8 & Drei Gleichungen & 75\\
9 & Wie weit reicht die Vorhersage & 83\\
\multicolumn{3}{@{}l}{\color{akzent}\bfseries\small IV.~Erkennen und Messen}\\[1pt]
10 & Ein Kanal genügt & 93\\
11 & Die Dimension des Attraktors & 101\\
12 & Der Test gegen den Zufall & 109\\
\multicolumn{3}{@{}l}{\color{akzent}\bfseries\small V.~Raum und Zeit}\\[1pt]
13 & Die Welle und die Spirale & 119\\
14 & Viele & 127\\
\multicolumn{3}{@{}l}{\color{akzent}\bfseries\small VI.~Stand}\\[1pt]
15 & Wo die Vorhersage endet & 137\\
\multicolumn{3}{@{}l}{\color{akzent}\bfseries\small VII.~Anhang}\\[1pt]
A & Formelsammlung und Übersicht & 145\\
B & Kommentierte Quellen & 157\\
\end{longtable}}

\section*{Band 29 --- Die Form der Lösungen}
\addcontentsline{toc}{section}{Band 29 --- Die Form der Lösungen}
\nopagebreak
{\small\itshape\color{grau} Algebraische Geometrie $\cdot$ 188 Seiten}\par\vspace{4pt}
\nopagebreak
{\small\renewcommand{\arraystretch}{1.28}
\begin{longtable}{@{}P{0.9cm}P{6.6cm}P{1.0cm}@{}}
\multicolumn{3}{@{}l}{\color{akzent}\bfseries\small I.~Nullstellen und Ideale}\\[1pt]
1 & Was eine Gleichung zeichnet & 7\\
2 & Die Gleichungen einer Varietät & 17\\
3 & Hilberts Nullstellensatz & 25\\
\multicolumn{3}{@{}l}{\color{akzent}\bfseries\small II.~Rechnen mit Idealen}\\[1pt]
4 & Ordnen und Teilen & 37\\
5 & Gröbner-Basen & 45\\
6 & Veränderliche loswerden & 55\\
\multicolumn{3}{@{}l}{\color{akzent}\bfseries\small III.~Die Gestalt}\\[1pt]
7 & Wie groß eine Varietät ist & 67\\
8 & Glatt und nicht glatt & 79\\
9 & Das Unendliche anschließen & 89\\
\multicolumn{3}{@{}l}{\color{akzent}\bfseries\small IV.~Kurven}\\[1pt]
10 & Die Gruppe auf der Kurve & 101\\
11 & Kurven, die man durchlaufen kann & 111\\
12 & Das Geschlecht & 121\\
\multicolumn{3}{@{}l}{\color{akzent}\bfseries\small V.~Anwendungen}\\[1pt]
13 & Modelle als Varietäten & 133\\
14 & Gelenke und Geheimnisse & 145\\
\multicolumn{3}{@{}l}{\color{akzent}\bfseries\small VI.~Stand}\\[1pt]
15 & Die Form der Lösungen & 159\\
\multicolumn{3}{@{}l}{\color{akzent}\bfseries\small VII.~Anhang}\\[1pt]
A & Formelsammlung und Übersicht & 167\\
B & Kommentierte Quellen & 179\\
\end{longtable}}

\section*{Band 30 --- Was die Symmetrie verrät}
\addcontentsline{toc}{section}{Band 30 --- Was die Symmetrie verrät}
\nopagebreak
{\small\itshape\color{grau} Darstellungstheorie $\cdot$ 196 Seiten}\par\vspace{4pt}
\nopagebreak
{\small\renewcommand{\arraystretch}{1.28}
\begin{longtable}{@{}P{0.9cm}P{6.6cm}P{1.0cm}@{}}
\multicolumn{3}{@{}l}{\color{akzent}\bfseries\small I.~Gruppen und ihre Darstellungen}\\[1pt]
1 & Symmetrie als Gruppe & 7\\
2 & Gruppen, die wirken & 17\\
3 & Die Bausteine & 27\\
\multicolumn{3}{@{}l}{\color{akzent}\bfseries\small II.~Charaktere}\\[1pt]
4 & Die Spur sagt alles & 39\\
5 & Tafeln bauen & 49\\
6 & Die Schwingungen des Moleküls & 61\\
\multicolumn{3}{@{}l}{\color{akzent}\bfseries\small III.~Konstruktionen}\\[1pt]
7 & Einschränken und Induzieren & 73\\
8 & Die symmetrische Gruppe & 83\\
9 & Produkte & 93\\
\multicolumn{3}{@{}l}{\color{akzent}\bfseries\small IV.~Kontinuierliche Gruppen}\\[1pt]
10 & Die Drehgruppe & 105\\
11 & Spin & 115\\
12 & Fourier auf Gruppen & 125\\
\multicolumn{3}{@{}l}{\color{akzent}\bfseries\small V.~Moleküle, Kristalle, Teilchen}\\[1pt]
13 & Benzol & 137\\
14 & Kristalle, Eisen, Quarks & 147\\
\multicolumn{3}{@{}l}{\color{akzent}\bfseries\small VI.~Stand}\\[1pt]
15 & Stand & 161\\
\multicolumn{3}{@{}l}{\color{akzent}\bfseries\small VII.~Anhang}\\[1pt]
A & Formelsammlung und Übersicht & 173\\
B & Kommentierte Quellen & 187\\
\end{longtable}}

\section*{Band B --- Die Architektur der Gewissheit}
\addcontentsline{toc}{section}{Band B --- Die Architektur der Gewissheit}
\nopagebreak
{\small\itshape\color{grau} Beweistheorie $\cdot$ 182 Seiten}\par\vspace{4pt}
\nopagebreak
{\small\renewcommand{\arraystretch}{1.28}
\begin{longtable}{@{}P{0.9cm}P{6.6cm}P{1.0cm}@{}}
\multicolumn{3}{@{}l}{\color{akzent}\bfseries\small I.~Die Frage}\\[1pt]
1 & Was ist ein Beweis? & 3\\
\multicolumn{3}{@{}l}{\color{akzent}\bfseries\small II.~Woher das alles kommt}\\[1pt]
2 & Rechnen ohne Beweisen: Die Welt vor Griechenland & 11\\
3 & Das griechische Wunder & 13\\
4 & Euklid: Die Erfindung des Systems & 17\\
5 & Archimedes: Finden und Beweisen sind zweierlei & 21\\
6 & Zwischenspiel: Algebra, Induktion, und die Rückkehr der Zahl & 23\\
7 & Die Krise der Infinitesimalien & 25\\
8 & Die Rigorisierung: Was ist eine reelle Zahl? & 27\\
9 & Nichteuklidische Geometrie: Der erste Unabhängigkeitsbeweis & 31\\
10 & Cantor: Das aktual Unendliche & 35\\
11 & Die Grundlagenkrise & 39\\
\multicolumn{3}{@{}l}{\color{akzent}\bfseries\small III.~Die Maschinerie}\\[1pt]
12 & Aussagenlogik: Das Skelett des Schließens & 45\\
13 & Prädikatenlogik: Die Sprache der Mathematik & 51\\
14 & Beweiskalküle: Wie man wirklich schließt & 57\\
15 & ZFC: Das Fundament, auf dem alles steht & 61\\
\multicolumn{3}{@{}l}{\color{akzent}\bfseries\small IV.~Der Werkzeugkasten}\\[1pt]
16 & Die elementaren Beweisformen & 69\\
17 & Induktion: Warum sie funktioniert & 73\\
18 & Diagonalisierung: Das mächtigste Muster & 77\\
19 & Zählen, Extremalprinzip, Invarianten & 79\\
20 & Nichtkonstruktive Methoden & 81\\
21 & Anatomie: Sechs Beweise unter dem Mikroskop & 85\\
\multicolumn{3}{@{}l}{\color{akzent}\bfseries\small V.~Die Grenzen}\\[1pt]
22 & Berechenbarkeit: Was eine Maschine kann & 93\\
23 & Gödels erster Unvollständigkeitssatz & 101\\
24 & Der zweite Unvollständigkeitssatz & 107\\
25 & Was Gödels Sätze nicht besagen & 109\\
26 & Unabhängigkeit: Die Kontinuumshypothese & 111\\
27 & Große Kardinalzahlen und die Suche nach neuen Axiomen & 115\\
28 & Konkrete Unabhängigkeit & 119\\
29 & Busy Beaver: Die Grenze in Zahlen & 123\\
30 & Beweistheorie: Die Vermessung der Systeme & 127\\
31 & Reverse Mathematik: Welche Axiome braucht ein Satz? & 129\\
32 & Beweiskomplexität: Wenn Beweise zu lang sind & 131\\
\multicolumn{3}{@{}l}{\color{akzent}\bfseries\small VI.~Andere Fundamente}\\[1pt]
33 & Konstruktive Mathematik & 135\\
34 & Typentheorie und Homotopietypentheorie & 137\\
35 & Kategorientheorie und Topoi & 139\\
\multicolumn{3}{@{}l}{\color{akzent}\bfseries\small VII.~Beweis heute}\\[1pt]
36 & Computerbeweise & 143\\
37 & Beweisassistenten und die Formalisierung der Mathematik & 145\\
38 & Künstliche Intelligenz und der Beweis & 147\\
39 & Die Soziologie und Epistemologie des Beweises & 151\\
\multicolumn{3}{@{}l}{\color{akzent}\bfseries\small VIII.~Epilog}\\[1pt]
40 & Wo sind die Grenzen wirklich? & 155\\
\multicolumn{3}{@{}l}{\color{akzent}\bfseries\small IX.~Anhänge}\\[1pt]
A & Symbolverzeichnis & 161\\
B & Zeitleiste & 163\\
C & Glossar & 167\\
D & Kommentierte Literatur & 169\\
\end{longtable}}

\section*{Gesamtumfang}
\addcontentsline{toc}{section}{Gesamtumfang}

\begin{merksatz}
\boxtitel{Zweiunddreißig Bände}
\begin{center}
{\small
\renewcommand{\arraystretch}{1.5}
\begin{tabular}{@{}P{3.4cm}P{2.0cm}P{2.4cm}@{}}
\toprule
\textbf{Abschnitt} & \textbf{Bände} & \textbf{Seiten}\\
\midrule
Erste Werkhälfte (0--15, 29, 30, B) & $19$ & $3348$\\
Zweite Werkhälfte (16--28) & $13$ & $2696$\\
\midrule
\textbf{Zusammen} & $\mathbf{32}$ & $\mathbf{6044}$\\
\bottomrule
\end{tabular}}
\end{center}

\begin{center}
\begin{minipage}{0.88\textwidth}
\kursiv{Der vorliegende Rahmenband ist nicht mitgezählt. \textbf{Die zweite Werkhälfte umfasst mit dreizehn Bänden knapp fünfundvierzig Prozent des Gesamtumfangs} --- zwölf wenden an, einer trägt eine Grundlage nach. Die Bände 29 und 30 sind Grundlagenbände, die nach der zweiten Hälfte geschrieben wurden und der ersten zugerechnet sind.}
\end{minipage}
\end{center}
\end{merksatz}